\documentclass[10pt,a4paper]{amsart}
\usepackage[margin=19mm]{geometry}
\usepackage{amsmath,amssymb,mathtools}
\usepackage[scr=rsfs,cal=euler]{mathalfa}
\usepackage{xfrac}
\usepackage[unicode,hidelinks,bookmarks=false]{hyperref}
\newcommand{\ie}{i.e.\ }
\newcommand{\cf}{cf.\ }
\newcommand{\resp}{respectively\ }
\newcommand{\Subsection}{\subsection}
\providecommand{\nobreakdash}{\nobreak\hbox{-}}
\setlength{\emergencystretch}{2em}
\multlinegap0pt
\usepackage{bm}
\usepackage{bbm}
\usepackage{dsfont}
\usepackage{xspace}
\usepackage{cancel}
\usepackage{mathtools}
\usepackage[shortlabels]{enumitem}
\usepackage{amsthm}
\makeatletter
\@ifundefined{theorem}{\newtheorem{theorem}{Theorem}}{}
\@ifundefined{lemma}{\newtheorem{lemma}[theorem]{Lemma}}{}
\makeatother
\title[Toward Minimum Graphic Parity Networks]{Toward Minimum Graphic Parity Networks}
\author{Yixin Cao, Yiren Lu, Junhong Nie, Xiaoming Sun, Guojing Tian}
\date{Source: arXiv:2509.10070v1 (CC BY 4.0)}
\begin{document}
\maketitle

\noindent\emph{Source and licence.} Toward Minimum Graphic Parity Networks. Version arXiv:2509.10070v1 (CC BY 4.0), distributed under the Creative Commons Attribution 4.0 International licence, as recorded in the arXiv licence metadata for this exact version and shown on the abstract page https://arxiv.org/abs/2509.10070v1. Full rights notice: licenses/arxiv-2509.10070-CC-BY-4.0.txt.\par\smallskip

\noindent\textit{Editorial scope.} Counted unit: the Lemma whose source label is \texttt{lem:analysis} in the article (source0.tex lines 637--643, source label \texttt{lem:analysis}), delivered together with the article's own complete original proof of it (source0.tex lines 644--730, one \texttt{proof} environment of 87 lines). No mathematics is written, repaired or re-derived: statement and proof are the article's own text line for line. Required context retained uncounted: none. Every definition, display and label of those lines is retained unchanged. Omitted from this package and not redistributed: 1--636, 731--1493. Nothing inside a retained excerpt is omitted or altered: every retained body is delivered exactly as the article's lines are written, so no omission receipt of any kind accompanies this package. Citations: the retained excerpts carry no citation command at all, so no bibliography entry of the article is transcribed and no reference is redistributed. Reference adaptation: none was needed and none was made; every reference inside the delivered unit resolves inside it, so no unresolved reference or citation is produced. Printed slips kept literal: no printed word, symbol, hypothesis or number of the counted statement or of its proof was altered. Layout and class: the article's own class is \texttt{article} with packages including fullpage, subcaption, makecell, amsthm, amsmath, amssymb, caption, complexity, enumerate, enumitem, microtype, booktabs, hyperref, cleveref; the wrapper is the lane's pinned \texttt{amsart} editorial wrapper at 10pt on A4, which restates the theorem-like environments the retained text needs and adds the article's own macro definitions that the retained text needs (none), with extra wrapper packages (bm, bbm, dsfont, xspace, cancel, mathtools, enumitem). The delivered numbering is the unit's own. Assets: no retained span contains a figure environment, a graphics-inclusion command, a PDF-page inclusion, a table or a TikZ picture; no image file is copied into this package, the package declares no asset dependency and no image file is redistributed with it. Licence: version arXiv:2509.10070v1 of this article is distributed under the Creative Commons Attribution 4.0 International licence, as recorded in the arXiv licence metadata for this exact version and shown on the abstract page https://arxiv.org/abs/2509.10070v1. A full rights notice is delivered at licenses/arxiv-2509.10070-CC-BY-4.0.txt. Characters: every retained excerpt is pure ASCII, so the delivered engine, XeLaTeX with its default Latin Modern text font, reports no missing character.
\begin{lemma}\label{lem:analysis}
  The expected number of line~12 being executed is upper bounded by
  \[
    4n+ \min_{1\le t\le n}\left\{2\sum_{i< t}d_i+\frac{(n-t)n\log n}{d_{t}}\right\},
  \]
  where $d_1,d_2,\ldots,d_n$ are the degrees of the vertices sorted in ascending order.
\end{lemma}

\begin{proof}
  Line~12 is only executed when~$K\ne \emptyset$ and~$i\not\in \mathrm{term}(j)$ at the beginning of the~$i$th iteration.  With permutation~$\pi$, the number of line~12 being executed is the number of pairs~$(u, v)$ such that there exists an extra vertex~$w$ satisfying
  \begin{enumerate}[label = (\roman*),left=10pt]
    \item $u w, v w\in E(G)$;
    \item $u$ is not the last neighbor of~$v$ in~$\pi$;
    \item $\pi(u) < \pi(v) < \pi(w)$; and
    \item $w x\not\in E(G)$ for all vertices~$x$ with $\pi(u) < \pi(x) < \pi(v)$.
  \end{enumerate}
  Therefore, it cannot be more than the number of pairs~$(u, v)$ such that there exists an extra vertex~$w$ merely satisfying (iv) and
  \begin{enumerate}[label = (\roman*),left=10pt]
    \setcounter{enumi}{3}
    \item $u w, v w\in E(G)$ and $\pi(u) < \pi(v)$.
  \end{enumerate}
  Now we obtain an upper bound on the number of such pairs.

  For any permutation~$\pi \in {\mathfrak {S}}_{n}$,
  let~$P(\pi)$ denote all such pairs when the algorithm is executed using permutation $\pi$, and let~$P(\pi, j)$ denote the subset of~$P(\pi)$ in which the first vertex is fixed by~$\pi(u) = j$.
  Then the expected number of such pairs is
  \begin{align}
    \notag
    \mathop{\mathbb{E}}_{\pi \in {\mathfrak {S}}_{n}} |P(\pi)| = & \mathop{\mathbb{E}}_{\pi \in {\mathfrak {S}}_{n}} \sum_{j = 1}^{n} |P(\pi, j)|
    \\      \notag
    =                                                            & \sum_{j = 1}^{n} \mathop{\mathbb{E}}_{\pi \in {\mathfrak {S}}_{n}}  |P(\pi, j)|
    \\      \notag
    =                                                            & \sum_{j=1}^n
    \mathop{\mathbb{E}}_{X\in {[n]\choose j}}
    \mathop{\mathbb{E}}_{\substack{{\pi \in {\mathfrak {S}}_{X}}                                                                                           \\  \forall x\in X. \pi(x) > n - j}}
    |P(\pi, n - j + 1)|
    \\
    \le                                                          & \sum_{j=1}^n \mathop{\mathbb{E}}_{\pi \in {\mathfrak {S}}_{n}} |P(\pi, 1)|.\label{eq:4}
  \end{align}
  Therefore, it reduces to bounding the expected size of~$P(\pi, 1)$.
  Consider any fixed~$t\in [n]$.
  If~$d(\pi^{-1}(1)) \le d_{t}$, then
  \begin{equation}
    \label{eq:5}
    |P(\pi, 1)| \le d_{t}.
  \end{equation}
  We now consider the nontrivial case, where~$d(\pi^{-1}(1)) > d_{t}$.
  Note that~$|P(\pi, 1)|$ is the number of vertices~$v$ such that there exists another vertex~$w$ whose first two neighbors in~$\pi$ are~$\pi^{-1}(1)$ and~$v$; we say that it is \emph{witnessed by~$w$}.
  For each vertex~$w\in N(\pi^{-1}(1))$, let~$X_{w}$ be the index of the second neighbor of~$x$ in~$\pi$.  Then
  \[
    \mathbb{P}\left( X_{w}=i\right) = \left(1-\frac {d(w)-1}{n-2}\right)\cdots \left(1-\frac {d(w)-1}{n-(i-1)}\right) \frac{d(w)-1}{n-i} >                                                   \left(1- \frac {d(w)-1}{n} \right)^{i - 2} \frac{d(w)-1}{n}.
  \]
  Letting~$\alpha = 1-\frac {d(w)-1}{n}$, we have
  \begin{align}
    \notag
      & \mathbb{P}\left(X_{w}\le {\frac {n \log n}{d(w)-1}} \right)
    \\
    \notag
    = & \mathbb{P}\left(X_{w}=2  \right) + \cdots +
    \mathbb{P}\left(X_{w}=\left\lfloor {\frac {n \log n}{d(w)-1}} \right\rfloor \right)
    \notag
    \\
    > & \frac{d(w)-1}{n} + \cdots + \alpha^{\left\lfloor {\frac {n \log n}{d(w)-1}} \right\rfloor - 2} \left(\frac{d(w)-1}{n}\right)
    \notag
    \\
    = & \left(\frac{1 - \alpha^{\left\lfloor {\frac {n \log n}{d(w)-1}} \right\rfloor - 1}}{1 - \alpha}\right) \left(\frac{d(w)-1}{n}\right)
    \notag
    \\
    = & {1 - \alpha^{\left\lfloor {\frac {n \log n}{d(w)-1}} \right\rfloor - 1}}.
    \notag
    \label{eq:2}
  \end{align}
  If~$d(w) > d_t$, then~$\frac {n\log n}{d(w)-1} \le \frac {n\log n}{d_t}$, and
  \begin{align}
    \mathbb{P}\left(X_{w}>\frac {n\log n}{d_t}\right)< \alpha^{\left\lfloor {\frac {n \log n}{d(w)-1}} \right\rfloor - 1} =  \left(1-\frac {d(w)-1}{n}\right)^{\left\lfloor {\frac {n \log n}{d(w)-1}} \right\rfloor - 1} < \frac{4}{n}.
  \end{align}
  Since the number of such vertices~$w$ is less than $n-1$ , then all such vertices witness at most~$\frac {n\log n}{d_t} + 4$ vertices.
  On the other hand, each vertex~$w$ with~$d(w)\le d_t$ witnesses at most one vertex.
  In summary, for each~$u\in V(G)$,
  \begin{equation}
    \label{eq:3}
    \mathop{\mathbb{E}}_{\substack{\pi \in {\mathfrak {S}_{n}}\\ \pi(u) = 1}} |P(\pi, 1)| <
    \frac {n\log n}{d_t} + 4 +
    \sum_{\substack{{w\in N(u)}\\ {d(w)\le d_t}}}1.
  \end{equation}

  Combining~\eqref{eq:4},~\eqref{eq:5}, and~\eqref{eq:3}, we conclude that the expected number of line~12 being executed is less than
  \begin{equation}
    \label{eq:6}
    \sum_{d_i\le t} d_i + \sum_{\substack{u \in V(G)\\ d(u) > d_{t}}} \left(\frac {n\log n}{d_t} + 4 +
    \sum_{\substack{{w\in N(u)}\\ {d(w)\le d_t}}}1\right)
    = 2\sum_{d_i\le t} d_i + \sum_{d_i>t} \frac {n\log n}{d_t} + 4n.
  \end{equation}
  The statement follows because~\eqref{eq:6} holds for all~$t\in [n]$.
\end{proof}

\end{document}
